\documentclass[11pt,a4paper]{article}
\usepackage[margin=2.5cm]{geometry}
\usepackage{fontspec}
\setmainfont{Times New Roman}
\usepackage{amsmath,amssymb}
\usepackage{booktabs,tabularx,array}
\usepackage{enumitem}
\usepackage{xcolor}
\usepackage{float}
\usepackage{titlesec}
\usepackage[hidelinks]{hyperref}

\titleformat{\section}{\Large\bfseries\color{black}}{\thesection}{0.6em}{}
\titleformat{\subsection}{\large\bfseries\color{black}}{\thesubsection}{0.6em}{}
\titleformat{\subsubsection}{\normalsize\bfseries\color{black}}{\thesubsubsection}{0.6em}{}

\setlength{\parindent}{0pt}
\setlength{\parskip}{5pt}
\linespread{1.25}
\renewcommand{\arraystretch}{1.15}
\renewcommand{\tabularxcolumn}[1]{>{\raggedright\arraybackslash}p{#1}}
\setlength{\emergencystretch}{2em}

\newcommand{\Cp}{\ensuremath{C_p}}
\newcommand{\CD}{\ensuremath{C_D}}
\newcommand{\ReD}{\ensuremath{\mathit{Re}}}

\pagestyle{plain}
\begin{document}

\begin{center}
{\LARGE\bfseries Cylinder Flow: Uncertainty and Error Propagation}\\[4pt]
{\large Companion notes to \texttt{cylinder\_lab.m}, \texttt{cpUncertainty.m}, \texttt{runUncertainty.m}}\\[8pt]
{\small Engineering Practice 2 -- Aerodynamics 2 \quad\textbar\quad RODH 115 (ex L50)
\quad\textbar\quad $d = 102$\,mm, $b = 460$\,mm, $d/b = 0.2217$}
\end{center}

\vspace{4pt}

\section{Purpose and conventions}

\subsection{Division of labour between the three files}

The calculation is split into three files so that each formula has exactly one home: change an
assumption once and nothing goes stale. \texttt{cpUncertainty.m} owns everything whose independent
variable is \Cp{} --- $u(\Cp)$ per tapping, $u(\CD')$, and the angular part of $u(\theta_{\text{sep}})$.
\texttt{runUncertainty.m} owns everything else: the dynamic-head-to-velocity-to-Reynolds chain, plus
the tapping-angle errors. \texttt{cylinder\_lab.m} only calls, prints and plots; it derives nothing.

Both modules are optional. The main script checks \texttt{exist} before calling, and falls back to
NaN with a printed note. The point is that forgetting to copy one file should not stop a lab session.

\subsection{Symbols, coverage factor, presentation}

Every $u$ here is a $1\sigma$ standard uncertainty, not a 95\,\% interval. When a 95\,\% figure is
needed, $k = 2$ is used (\texttt{cfg.sepUncK}) --- an engineering convention, since this experiment
does not carry enough degrees of freedom to fix the coverage factor properly.

Error bars on all figures are $1\sigma$, stated in the titles. In the results CSV, columns starting
\texttt{u\_} are $1\sigma$ and columns starting \texttt{U95} are $k=2$. Keeping both lets you explain
in the viva the difference between the bar you drew and the number you quoted.

\section{Measurement chain and inputs}

\subsection{From column height to \Cp}

Each tapping gives a height $h_i$, static gives $h_s$, total gives $h_t$:

\begin{equation}
\Cp{}_i = \frac{h_i - h_s}{h_t - h_s} = \frac{A}{B}
\end{equation}

Ethanol density, gravity, the scale zero and whether the tubes connect at the top or bottom of the
board all cancel in the ratio. That is also why the connection direction does not change \Cp:
reversed connection flips every column, so numerator and denominator flip together. The one thing
that genuinely needs a physical decision at the rig is which Pitot tube is total. The criterion is
that hole 2 ($\theta = 0$, stagnation) has surface pressure equal to total pressure, so the tube
reading equal to $h(\text{hole } 2)$ is the total one. If both are far from stagnation the script
takes the larger and prints a note.

\subsection{From column height to $U$ and \ReD}

\begin{equation}
q = \rho_e g\,|\Delta h|, \qquad U = \sqrt{\frac{2q}{\rho}}, \qquad \ReD = \frac{\rho U d}{\mu}
\end{equation}

\begin{equation}
\rho = \frac{p}{RT}, \qquad \mu = \mu_0\left(\frac{T}{T_0}\right)^{1.5}\frac{T_0 + S}{T + S}
\end{equation}

Here $\Delta h = h_t - h_s$, the same $B$ as above. Note that $\rho$ appears both in $\rho U d$ and
inside $U = \sqrt{2q/\rho}$; that duplication fixes a coefficient of $+0.5$ in Section~\ref{sec:re}
and is the easiest sign to get wrong.

\subsection{Inputs and assumptions}

\begin{table}[H]
\centering
\caption{Every input to the uncertainty model}
\small
\begin{tabularx}{\textwidth}{@{}l p{3.4cm} p{2.5cm} c X@{}}
\toprule
Symbol & Meaning & Value & Type & Source / note \\
\midrule
$u$ & uncertainty of one column reading & 0.5 mm & A+B merged, conservative &
\texttt{cfg.uRead}; fallback when no repeats \\
$u_{\text{scale}}$ & type B from the scale graduation & 1.0 mm $\to$ 0.289 mm & B &
\texttt{cfg.uScale}, rectangular $(1/2)/\sqrt{3}$ \\
$u_\theta$ & standard uncertainty of one tapping angle & 0.29$^\circ$ & B &
\texttt{cfg.uTheta}, from a $\pm0.5^\circ$ tolerance \\
$u_T$ & air temperature & 0.3 K & B & \texttt{cfg.uTair} \\
$u_p$ & barometric pressure & 50 Pa & B & \texttt{cfg.uPbaro} \\
$n$ & repeats per reading & 3 & --- & three columns per row in \texttt{readings.csv} \\
$k$ & coverage factor for the 95\,\% figure & 2 & --- & \texttt{cfg.sepUncK} \\
\bottomrule
\end{tabularx}
\end{table}

All entries except $u_{\text{scale}}$ are quantities whose true value is unknown and only a range is
known, so the default GUM treatment applies: rectangular distribution, standard uncertainty equal
to half-width divided by $\sqrt{3}$. The $0.5^\circ$ behind $u_\theta$ is an assumption
(Section~\ref{sec:theta}), not data from the handout.

\section{Uncertainty of \Cp}

\subsection{The three partial derivatives}

With $A = h_i - h_s$, $B = h_t - h_s$ and the three readings independent:

\begin{equation}
\frac{\partial \Cp}{\partial h_i} = \frac{1}{B}, \qquad
\frac{\partial \Cp}{\partial h_t} = -\frac{\Cp}{B}, \qquad
\frac{\partial \Cp}{\partial h_s} = \frac{\Cp - 1}{B}
\end{equation}

The second and third are not constants: $h_t$ and $h_s$ sit in both numerator and denominator, so
their weight depends on \Cp{} itself. This is why $u(\Cp)$ is larger at the suction peak than at
stagnation --- there $\Cp \approx -2.3$ and the $(\Cp - 1)$ factor amplifies the static-tube error
more than threefold.

\subsection{Single-point combination}

\begin{equation}
u(\Cp) = \frac{u}{B}\sqrt{1 + \Cp^2 + (\Cp - 1)^2}
\end{equation}

$B = |h_t - h_s|$ in mm, so the right-hand side is dimensionless. At low speed $B \approx 18.7$\,mm
and at high speed $B \approx 72.4$\,mm: the entire low-versus-high difference comes from $B$.

\subsection{Correlation between tappings in one run}

The tappings are not independent --- they share one $h_s$ and one $h_t$:

\begin{equation}
\operatorname{Cov}(C_{p,i}, C_{p,j}) = \left(\frac{u}{B}\right)^2
\bigl[\,C_{p,i} C_{p,j} + (C_{p,i} - 1)(C_{p,j} - 1)\,\bigr], \qquad i \neq j
\end{equation}

The first term comes from the shared total tube, the second from the shared static tube, and both
carry a positive sign: a static tube reading high pushes every \Cp{} down together, so the errors
add rather than cancel. Dropping this $28\times28$ covariance when integrating for \CD{} understates
$u(\CD)$ by roughly a factor of two.

\subsection{Numbers}

\begin{table}[H]
\centering
\caption{Single-point $u(\Cp)$ ($B = 18.7$\,mm low, $72.4$\,mm high; $u = 0.5$\,mm)}
\small
\begin{tabular}{@{}l c c c c@{}}
\toprule
Case & $\Cp = +1$ (stagnation) & $\Cp = -2.3$ (peak) & $\Cp = 0$ & $u(\CD)$ \\
\midrule
low speed, 15.5 m/s & 0.038 & 0.110 & 0.055 & 0.040 \\
high speed, 30.5 m/s & 0.010 & 0.028 & 0.014 & 0.010 \\
low speed, 1 mm reading error & 0.081 & --- & --- & 0.081 \\
\bottomrule
\end{tabular}
\end{table}

The 1\,mm row is the range estimate: reading the column to 1\,mm instead of 0.5\,mm doubles
$u(\CD)$ to 0.081. That is why coarse estimation at low speed is not acceptable.

\section{Uncertainty of \CD}

\subsection{Where the weights come from}

$\CD = \int_0^{\pi} \Cp \cos\theta \,\mathrm{d}\theta$, evaluated with \texttt{trapz} rather than
uniform summation, because the 28 angles are not equally spaced: $5^\circ$ steps near
$0$--$20^\circ$, $5$--$10^\circ$ in the middle, $10$--$15^\circ$ towards $180^\circ$. \texttt{trapz}
assigns interior weights $(\theta_{i+1} - \theta_{i-1})/2$ and half that at the two ends. These
weights set both the value of \CD{} and the coefficients $g_i = w_i \cos\theta_i$ used in propagation.

\subsection{Closed form}

Let $s = u/B$ and substitute the three partials into $\CD = \sum g_i C_{p,i}$. Expanding the covariance
fully, the cross terms collapse into

\begin{equation}
u(\CD) = s\sqrt{\sum_i g_i^2 + \Bigl(\sum_i g_i (C_{p,i} - 1)\Bigr)^2 + \Bigl(\sum_i g_i C_{p,i}\Bigr)^2}
\end{equation}

Reading the three terms: the first is each tapping's own random error (uncorrelated, so plain
sum of squares); the second is the static-tube common-mode error with weight $(C_{p,i} - 1)$; the third
is the total-tube common-mode error with weight $C_{p,i}$. Both common-mode terms are squared after
summation, which is precisely what a correlated error looks like.

\subsection{Per-tapping measured version}

With three repeats per tapping in \texttt{readings.csv}, $u$ becomes a per-tapping
$u_{\text{tap},i}$ and the expression splits:

\begin{align}
A &= \frac{1}{B^2}\sum_i g_i^2\Bigl[\,u_{\text{tap},i}^2 + (C_{p,i} u_t)^2 + \bigl((C_{p,i}-1)u_s\bigr)^2\,\Bigr] \\
B_1 &= \left(\frac{u_t}{B}\right)^2\Bigl[\Bigl(\sum_i g_i C_{p,i}\Bigr)^2 - \sum_i (g_i C_{p,i})^2\Bigr] \\
B_2 &= \left(\frac{u_s}{B}\right)^2\Bigl[\Bigl(\sum_i g_i (C_{p,i}-1)\Bigr)^2 - \sum_i \bigl(g_i (C_{p,i}-1)\bigr)^2\Bigr] \\
u(\CD) &= \sqrt{A + B_1 + B_2}
\end{align}

The subtracted terms in $B_1$ and $B_2$ are the within-tapping parts already counted in $A$;
without them the double counting would inflate the result.

\subsection{The two angular terms}

Tapping positions were never in the model above. Position errors come in two flavours with
completely different propagation; they cannot be merged into one number.

\textbf{Per-tapping random part.} Tapping $i$ sits at $\theta_i + e_i$ with $e_i$ independent and
standard deviation $u_\theta$. Displace that tapping alone by $\pm u_\theta$, re-integrate the whole
\CD, take the central difference, and combine over all 28:

\begin{equation}
u(\CD)_{\text{taps}} = \sqrt{\sum_i \left[\frac{\CD(\theta + e_i) - \CD(\theta - e_i)}{2}\right]^2}
\end{equation}

Re-integrating rather than writing an analytic derivative avoids a three-branch formula: moving one
tapping also changes the width of the two trapezoids on either side of it. On the demo data this
term evaluates to 0.0088.

\textbf{Rigid rotation.} Hole 1 ($-30^\circ$) and hole 7 ($+30^\circ$) are mirror images and should
read the same. A difference $\Delta C_{p,30}$ can only mean the whole ring sits at
$\Delta\theta$ relative to the flow. With $g$ the measured slope at $30^\circ$,
$\Cp(\mp 30 + \Delta\theta) \approx P \mp g\Delta\theta$, so subtracting gives
$\Delta C_{p,30} = -2g\Delta\theta$ and

\begin{equation}
\Delta\theta = -\frac{\Delta C_{p,30}}{2g}, \qquad
g = \left.\frac{\mathrm{d}\Cp}{\mathrm{d}\theta}\right|_{30^\circ}
\text{ from the $20^\circ$ and $40^\circ$ holes}
\end{equation}

Its effect on \CD{} is the difference obtained by shifting the whole curve,
$|\CD(\theta + \Delta\theta) - \CD(\theta)|$, treated as a bias and never averaged away against the
per-tapping term. On the demo data $\Delta\theta = -0.05^\circ$ and this term is 0.0022.

\subsection{Combined result}

\begin{table}[H]
\centering
\caption{Breakdown of $u(\CD')$ on the demo data (conservative 0.5\,mm model)}
\small
\begin{tabular}{@{}l c c c c c c@{}}
\toprule
run & $\CD'$ & reading & per-tap angle & rotation & total $u(\CD')$ & relative \\
\midrule
smooth, low & 1.109 & 0.040 & 0.0088 & 0.0022 & 0.0449 & 4.0\,\% \\
smooth, high & 0.925 & 0.010 & 0.0079 & 0.0022 & 0.0127 & 1.4\,\% \\
tripped, low & 0.465 & 0.022 & 0.0066 & 0.0048 & 0.0236 & 5.1\,\% \\
tripped, high & 0.597 & 0.007 & 0.0070 & 0.0009 & 0.0098 & 1.6\,\% \\
\bottomrule
\end{tabular}
\end{table}

Worth noticing: at high speed the angular term is nearly the size of the reading term (0.007 against
0.007). The pressure measurement has become good enough that the bottleneck is now geometry. That
belongs in the report, because it changes which dataset you should trust.

\section{Uncertainty of $\theta$}
\label{sec:theta}

\subsection{The 0.29$^\circ$ assumption}

Tapping position error is a machining tolerance with neither repeats nor a known distribution, so
it is treated as rectangular over $\pm0.5^\circ$:

\begin{equation}
u_\theta = \frac{0.5^\circ}{\sqrt{3}} = 0.289^\circ \approx 0.29^\circ
\end{equation}

The $\pm0.5^\circ$ is an assumption about dividing-head practice; the handout does not quote a
tolerance, which makes this the only number in the model set by judgement. If a real tolerance
$T$ becomes available, set \texttt{cfg.uTheta} to $T/\sqrt{3}$. All 28 holes share one value: they
come from one machining pass, are identically distributed, and there is no reason to make them differ.

\subsection{Uncertainty of the separation angle itself}

The separation point is found by a slope criterion: after the suction peak, the first place where
the slope is near zero and stays near zero over the following $12^\circ$. It is noise-sensitive,
estimated as slope noise over curvature:

\begin{equation}
u(\theta_{\text{sep}})_{\text{slope}} = \frac{\sigma_{\text{slope}}}{|\mathrm{d}^2 \Cp/\mathrm{d}\theta^2|}
\end{equation}

clipped below at half a tapping pitch (two holes cannot resolve a position between them) and above
at $45^\circ$. The upper clip fires when the separation point approaches $180^\circ$ and the tail
is nearly flat --- the criterion has effectively stopped working. The demo values are
$2.5^\circ$, $2.5^\circ$, $19.9^\circ$, $4.0^\circ$; the third is direct evidence of that failure and
should not be read as a well-measured angle.

\subsection{Combining}

\begin{equation}
u(\theta_{\text{sep}}) = \sqrt{u(\theta_{\text{sep}})^2_{\text{slope}} + u_\theta^2 + \Delta\theta^2}
\end{equation}

In practice $u_\theta = 0.29^\circ$ and $\Delta\theta \approx 0.1^\circ$ are far smaller than the
slope term, so the value barely changes; their role is to provide a floor when the data happen to
smooth into a very sharp criterion.

\section{Uncertainty of \ReD}
\label{sec:re}

\subsection{The chain and the $+0.5$}

Take logarithmic differentials, remembering that $\rho$ appears twice --- once in $\rho U d$ and
once through $U = \sqrt{2q/\rho}$:

\begin{align}
\mathrm{d}\ln U &= 0.5\,\mathrm{d}\ln q - 0.5\,\mathrm{d}\ln\rho \\
\mathrm{d}\ln \ReD &= 0.5\,\mathrm{d}\ln\rho + 0.5\,\mathrm{d}\ln q - \mathrm{d}\ln\mu
\end{align}

Hence $+0.5$ on $\rho$, not a minus sign. The three inputs are independent:

\begin{equation}
\frac{u(\ReD)}{\ReD} = \sqrt{\left(0.5\frac{u(q)}{q}\right)^2
+ \left(0.5\frac{u(\rho)}{\rho}\right)^2 + \left(\frac{u(\mu)}{\mu}\right)^2}
\end{equation}

with
$\dfrac{u(q)}{q} = \dfrac{\sqrt{u_s^2 + u_t^2}}{B}$,
$\dfrac{u(\rho)}{\rho} = \sqrt{\left(\dfrac{u_p}{p}\right)^2 + \left(\dfrac{u_T}{T}\right)^2}$,
$\dfrac{u(\mu)}{\mu} = \left|\dfrac{1.5}{T} - \dfrac{1}{T+S}\right|u_T$,
the last being the analytic sensitivity of Sutherland's law.

\subsection{The corrected \ReD{} carries one more term}

$V/V' = 1 + 0.25\,\CD'(d/b) + 0.82(d/b)^2$ depends on $\CD'$, which has its own uncertainty:

\begin{equation}
u(V/V') = 0.25\,(d/b)\,u(\CD')
\end{equation}

\begin{equation}
\frac{u(\ReD')}{\ReD'} = \sqrt{\left(\frac{u(\ReD)}{\ReD}\right)^2
+ \left(\frac{u(V/V')}{V/V'}\right)^2}
\end{equation}

On the demo this is only 0.22\,\% (low) and 0.05\,\% (high), because $0.25(d/b) \approx 0.055$
multiplied by $u(\CD') \approx 0.045$ is small. It is kept so that the correction itself is
accountable.

\subsection{Numbers}

\begin{table}[H]
\centering
\caption{Breakdown of $u(\ReD)$ (demo, conservative 0.5\,mm model)}
\small
\begin{tabular}{@{}l c c c c c c@{}}
\toprule
case & $u(q)/q$ & $u(\rho)/\rho$ & $u(\mu)/\mu$ & $u(\ReD')/\ReD'$ & $V/V'$ & $u(\ReD)/\ReD$ \\
\midrule
low speed & 3.79\,\% & 0.113\,\% & 0.079\,\% & 1.90\,\% & 0.221\,\% & 1.91\,\% \\
high speed & 0.98\,\% & 0.113\,\% & 0.079\,\% & 0.50\,\% & 0.049\,\% & 0.50\,\% \\
\bottomrule
\end{tabular}
\end{table}

The halved $q$ term still dominates; $\rho$ and $\mu$ together stay under 0.14\,\%. Conclusion:
$u(\ReD)$ can only be reduced by reading the column better, no matter how precisely the room
temperature is measured.

\section{Presentation rules}

\begin{table}[H]
\centering
\caption{How each quantity appears in the figures and the CSV}
\small
\begin{tabularx}{\textwidth}{@{}l X l X@{}}
\toprule
Quantity & figure & CSV column & how to quote it \\
\midrule
$u(\Cp)$ & nested bars per tapping: $1\sigma$ solid, $2\sigma$ thin pale & \texttt{u\_Cp\_peak} &
one value at the suction peak characterises the run \\
$u(\CD')$ & nested vertical bars on the \CD--\ReD{} plot & \texttt{u\_CD} & quote absolute and relative \\
$u(\ReD)$ & horizontal bar on the \CD--\ReD{} plot ($1\sigma$ only) & not written & percentage in figure and text only \\
tapping angle & horizontal bar on the \Cp--$\theta$ plot & --- & state it is an assumption \\
$u(\theta_{\text{sep}})$ & horizontal segment at the separation marker & \texttt{u\_thetaSep} &
always mention the criterion's own limitation \\
\bottomrule
\end{tabularx}
\end{table}

One rule: every point carries two levels --- $1\sigma$ as a solid capped bar, $2\sigma$ as a thin pale
bar drawn outside it --- so the reader never has to guess which sigma a bar represents. A 95\,\%
figure must still state $k=2$, and the bar length must not be quoted as $U_{95}$ ($2\sigma$ equals
$U_{95}$ only for a normal distribution with enough degrees of freedom).

Each figure in \texttt{results/} is written twice: \texttt{.fig} for editing in MATLAB and
\texttt{.pdf} (vector) for \texttt{\textbackslash includegraphics} in this report.
\texttt{cfg.figFormat} switches the export to \texttt{'eps'} or \texttt{'png'} if needed.

\section{Known limitations}

\begin{table}[H]
\centering
\caption{Limitations and their effect on the conclusions}
\small
\begin{tabularx}{\textwidth}{@{}p{4.0cm} X X@{}}
\toprule
limitation & effect & how to phrase it \\
\midrule
blockage ratio 22.2\,\%, Roshko calibrated to 13.6\,\%, 1.63$\times$ extrapolation &
the validity of the correction is itself in doubt; a few points of the 10\,\% correction may be extrapolation error &
report corrected and uncorrected side by side, never the corrected one alone \\
$\pm0.5^\circ$ tolerance is assumed &
the angular part of $u(\CD)$ may be off by a factor of two &
say ``assuming $\pm0.5^\circ$'', and give the value for $0.1^\circ$ \\
the $\pm30$ difference collapses into one equivalent rotation &
true asymmetry (hole depth, cylinder not perpendicular) is folded in &
call it ``equivalent rotation'', not ``alignment error'' \\
the separation criterion fails at the tail &
$u(\theta_{\text{sep}})$ reaches $19.9^\circ$ on tripped low (45$^\circ$ with the measured model) &
quote the suction-peak angle as well; it does not depend on the criterion \\
no cross-run correlation treatment &
same manometer, same tappings, so errors are partly common &
compare runs side by side, do not weight-average across them \\
\bottomrule
\end{tabularx}
\end{table}

\section{Code map}

When the viva asks where a number comes from, this table answers it. Line numbers refer to the
current version.

\begin{table}[H]
\centering
\caption{Where each output is computed}
\small
\begin{tabularx}{\textwidth}{@{}p{4.6cm} p{4.2cm} X@{}}
\toprule
output & file & implementation note \\
\midrule
\texttt{cfg.uRead / uScale / uTheta / uTair / uPbaro} & \texttt{cylinder\_lab.m}, lines 68--78 &
all assumptions in one block \\
$u(\Cp)$, $u(\CD')$, slope part of $u(\theta_{\text{sep}})$ & \texttt{cpUncertainty.m} &
covariance handled here and nowhere else \\
$\Delta\theta$ and $u(\CD)_{\text{taps}}$ & \texttt{runUncertainty.m}, lines 87--108 &
per-tapping central difference plus rigid shift \\
$u(\ReD)$ chain & \texttt{runUncertainty.m}, lines 60--80 &
$q \to U \to \ReD$ with $\rho$, $\mu$, $V/V'$ terms \\
$\Delta\theta$ from the $\pm30$ pair & \texttt{cylinder\_lab.m}, lines 289--304 &
$\Delta C_{p,30}$ with the $20^\circ/40^\circ$ central-difference slope \\
bar lengths & \texttt{sigmaCp / sigmaCD / thetaBar / separationBar} &
measured when available, conservative otherwise \\
CSV \texttt{u\_} columns & \texttt{writeCsv} & 29 columns, returned uncertainties only \\
\bottomrule
\end{tabularx}
\end{table}

\section{Viva questions}

\begin{enumerate}[leftmargin=1.6em,itemsep=4pt]
\item \textbf{Why is \Cp{} only a ratio of column differences?}\\
Ethanol density, gravity, scale zero and connection direction all cancel. Only the ratio of
differences matters. The sign convention does not change the result, but which Pitot tube is total
must be decided physically.

\item \textbf{Why is $u(\Cp)$ larger at the suction peak?}\\
$\partial \Cp/\partial h_s = (\Cp - 1)/B$; at $\Cp \approx -2.3$ this is more than three, so the
static-tube error is amplified.

\item \textbf{Why can't the per-tapping $u(\Cp)$ simply be added in quadrature for \CD?}\\
They share $h_s$ and $h_t$; the covariance is positive, and ignoring it roughly halves $u(\CD)$.

\item \textbf{Why \texttt{trapz} rather than uniform summation?}\\
The 28 angles are uneven ($5^\circ$, $10^\circ$, $15^\circ$ mixed); uniform weights would
over-weight the sparse region.

\item \textbf{Why are there two angular terms?}\\
Per-tapping errors dilute as $\sqrt{N}$ over 28 holes; a rigid rotation is common mode and does not
dilute at all, entering linearly.

\item \textbf{How does the $\pm30$ difference become an angle?}\\
By symmetry the two points should read the same; the difference means the ring is rotated.
$\Cp(\mp30 + \Delta\theta) \approx P \mp g\Delta\theta$ gives $\Delta C_{p,30} = -2g\Delta\theta$.

\item \textbf{Why measure $g$ instead of using the potential-flow $-4\sin 2\theta$?}\\
The boundary layer flattens \Cp{} near $30^\circ$; the measured slope is about 70\,\% of theory,
and using theory would systematically shrink $\Delta\theta$.

\item \textbf{Why is the coefficient on $\rho$ in $u(\ReD)$ equal to $+0.5$?}\\
$\rho$ enters through $\rho U d$ and through $U = \sqrt{2q/\rho}$; adding gives $+1 - 0.5 = +0.5$.

\item \textbf{What does 0.5 mm versus 1 mm cost?}\\
$u(\CD)$ goes from 0.040 to 0.081 --- it doubles. Reading accuracy sets the credibility of the whole
conclusion.

\item \textbf{Is the $19.9^\circ$ a data problem or a criterion problem?}\\
A criterion problem: the tail is nearly flat, the slope threshold sits at noise level and the
criterion stops working. Quote the suction-peak angle, which does not depend on it.
\end{enumerate}

\end{document}
